\documentclass[10pt]{article}
\usepackage[T1]{fontenc}
\usepackage[utf8]{inputenc}
\usepackage[french]{babel}
\usepackage{amsmath,amssymb,mathrsfs,mathpazo}
\usepackage{geometry}
\geometry{a4paper,margin=2cm}
\usepackage{graphicx,array,tabularx,tabularray,colortbl,multirow,multicol,booktabs,enumitem,pifont,ulem,fancybox,diagbox,cancel}
\usepackage{pstricks,pstricks-add,pst-plot,pst-tree,pst-node,pst-text,pst-func,pst-math,pst-bspline,pst-3dplot,pst-eucl,pst-solides3d}
\let\pscalculate\undefined
\usepackage{pst-bezier}
\usepackage{tikz,pgfplots,tkz-tab}
\usetikzlibrary{arrows,arrows.meta,patterns,calc,positioning,trees,decorations.markings,shapes.geometric,shapes.misc,shapes.arrows,intersections,decorations.pathmorphing,decorations.pathreplacing,backgrounds,shadows,babel}
\pgfplotsset{compat=1.18}
\usepgfplotslibrary{fillbetween}
\usepackage[np]{numprint}
\usepackage{listings,makecell,pas-tableur,twoopt}
\let\barre\undefined
\usepackage[pstricks]{tablvar}
\definecolor{darkgreen}{RGB}{0,100,0}
\providecommand{\red}{\color{red}}
\providecommand{\blue}{\color{blue}}
\providecommand{\green}{\color{green}}
\providecommand{\black}{\color{black}}
\providecommand{\white}{\color{white}}
\providecommand{\gray}{\color{gray}}
\providecommand{\cyan}{\color{cyan}}
\providecommand{\magenta}{\color{magenta}}
\providecommand{\yellow}{\color{yellow}}
\providecommand{\violet}{\color{violet}}
\providecommand{\orange}{\color{orange}}
\providecommand{\euro}{\texteuro}
\providecommand{\EUR}{\texteuro}
\providecommand{\ieme}{\textsuperscript{e}}
\providecommand{\degres}{\ensuremath{^\circ}}
\providecommand{\degree}{\ensuremath{^\circ}}
\providecommand{\vv}[1]{\overrightarrow{#1}}
\providecommand{\Coord}[1]{#1}
\providecommand{\decosix}{}
\providecommand{\starredbullet}{\ensuremath{\bullet}}
\pagestyle{empty}
\setlength{\parindent}{0pt}
\usepackage{adjustbox}\DeclareUnicodeCharacter{2236}{:}\DeclareUnicodeCharacter{2190}{\ensuremath{\leftarrow}}
\providecommand{\R}{\mathbb{R}}
\providecommand{\N}{\mathbb{N}}
\providecommand{\Z}{\mathbb{Z}}
\providecommand{\Q}{\mathbb{Q}}
\providecommand{\D}{\mathbb{D}}
\providecommand{\C}{\mathbb{C}}
\providecommand{\e}{{\mathrm{e}}}
\providecommand{\vect}[1]{\overrightarrow{#1}}
\providecommand{\vectt}[1]{\overrightarrow{\mathrm{#1}}}
\providecommand{\vvt}[1]{\overrightarrow{\mathrm{#1}}}
\providecommand{\np}[2][]{{#2~#1}}
\providecommand{\vv}[1]{{\overrightarrow{#1}}}
\renewcommand{\d}{{\mathrm{d}}}
\providecommand{\degres}{{^\circ}}
\providecommand{\degre}{{^\circ}}
\providecommand{\blue}{\color{blue}}
\providecommand{\red}{\color{red}}
\providecommand{\enskip}{\quad}
\providecommand{\llbracket}{\lbrack\!\lbrack}
\providecommand{\rrbracket}{\rbrack\!\rbrack}
\providecommand{\euro}{{\text{€}}}
\providecommand{\no}{{\text{n°}}}
\providecommand{\interval}[2]{{[#1;#2]}}
\providecommand{\intervFO}[2]{{[#1;#2[}}
\providecommand{\dotfill}{\ldots}

\providecommand{\ds}{\displaystyle}
\providecommand{\textding}[1]{#1}
\providecommand{\Oij}{\(\mathrm{O};\vec{\imath},\vec{\jmath}\)}
\providecommand{\Oijk}{\(\mathrm{O};\vec{\imath},\vec{\jmath},\vec{k}\)}
\providecommand{\Ouv}{\(\mathrm{O};\vec{u},\vec{v}\)}
\providecommand{\intervFF}[2]{\ensuremath{[#1;#2]}}
\renewcommand{\euro}{\ifmmode\text{\texteuro}\else\texteuro\fi}
\DeclareUnicodeCharacter{00B0}{\ifmmode^{\circ}\else\textdegree\fi}
\DeclareUnicodeCharacter{20AC}{\ifmmode\text{\texteuro}\else\texteuro\fi}
% Énoncés : Éducation nationale. Transcription et corrigés : APMEP, Denis Vergès et contributeurs. Découpage et classement : CoopMaths, Sébastien Lozano, Éric Elter et contributeurs. Adaptation de la présentation : Maths Explorer.
% CC BY-SA : https://coopmaths.fr/www/about/
% https://coopmaths.fr/alea/static/bac/2026/tex/bac_2026_06_sujet1_metropole_4.tex
\begin{document}
\noindent\textit{APMEP / CoopMaths. CC BY-SA. Présentation adaptée par Maths Explorer.}\par\medskip
On considère la fonction $f$ définie sur l'intervalle $]-1 ;+\infty[$ par :
	\[f(x)=a+\dfrac{b \ln (x+1)}{x+1}.\]
	où $a$ et $b$ sont des réels, et ln désigne la fonction logarithme népérien.

	On admet que la fonction $f$ est deux fois dérivable sur $]-1 ; +\infty [$.

	On note $f'$ sa fonction dérivée et $f''$ sa fonction dérivée seconde.

	\subsection*{Partie A}
	L'objectif de cette partie est de déterminer les valeurs de $a$ et $b$ de sorte que la courbe représentative de la fonction $f$ sur $]-1 ;+\infty[$ 
	corresponde à celle tracée dans le repère orthonormé ci-dessous :

	\begin{center}
		
		\begin{tikzpicture}
			\begin{axis}[/pgf/number format/.cd, use comma,
				x={12mm}, y={12mm},
				xmin = -0.5, xmax = 7.2,
				ymin = -0.8, ymax = 5.8,
				xtick distance=1, ytick distance=1,
				xticklabels={},yticklabels={},
				minor tick num = 0,
				grid=major, axis lines =center]
				\addplot [line width=1.2pt,color=red,smooth,samples=100,domain= -0.5:7.2 ]{1+(4*ln(x+1))/(x+1)};
				\node[below left] at (axis cs: 0,0) {0};
				\node[below] at (axis cs: 1,0) {1};
				\node[left] at (axis cs: 0,1) {1};
				\draw[shift={(axis cs: 0,1)}] (2pt,2pt)--(-2pt,-2pt) (-2pt,2pt)--(2pt,-2pt) node[below right]{A};
				\draw[dashed,line width=0.8pt] (axis cs: -1,-3)--(axis cs: 2,9) node[pos=0.55,left] {$T_{\mathrm{A}}$} ;
			\end{axis}
		\end{tikzpicture}













	\end{center}


	On précise que la droite $T_{\mathrm{A}}$ est la tangente à la courbe représentative de la fonction $f$ au point A$(0 ; 1)$.

	\begin{enumerate}
		\item Justifier que $a=1$.

		\item En utilisant le graphique :
		\begin{enumerate}
			\item Donner la valeur de $f'(0)$.

			Justifier.
			\item Donner le signe de $f''(1)$.

			Justifier.
		\end{enumerate}

		\item \begin{enumerate}
			\item Montrer que pour tout réel $x$ appartenant à $]-1 ;+\infty[$, on a :
			\[f'(x)=\dfrac{b(1-\ln (x+1))}{(x+1)^{2}}\]

			\item En déduire la valeur de $b$.
		\end{enumerate}
	\end{enumerate}


	\subsection*{Partie B}
	On admet dans la suite de l'exercice que la fonction $f$ est définie sur $]-1 ;+\infty[$ par :
	\[f(x)=1+\dfrac{4 \ln (x+1)}{x+1}\]

	\begin{enumerate}
		\item Justifier que la droite d'équation \quad $y=1$ \quad est asymptote à la courbe représentative de la fonction~$f$.

		\item Résoudre l'inéquation \quad $1-\ln (x+1)>0$ \quad sur $]-1 ;+\infty[$.

		\item On admet que \quad $\lim\limits_{x \to -1} f(x)=-\infty$.

		Dresser le tableau de variation complet de la fonction $f$ en indiquant la valeur exacte de son extremum.

		Justifier.

		\item Montrer que l'équation \quad $f(x)=1,5$ \quad admet une unique solution dans l'intervalle $[2 ; +\infty[$.

		En donner une valeur arrondie à $10^{-1}$.

		\item \begin{enumerate}
			\item Montrer que :
			\[\int_{0}^{2} \dfrac{\ln (x+1)}{x+1} \d x=\dfrac{1}{2}\big(\ln 3\big)^{2}.\]

			\item En déduire l'aire, en unité d'aire, du domaine délimité par l'axe des abscisses, la courbe représentative de la fonction $f$, l'axe des ordonnées et la droite d'équation \quad $x=2$.
		\end{enumerate}
	\end{enumerate}

\end{document}
